%% ===================================================================
%% KAPITOLA 6: ZÁVĚR
%% ===================================================================
\begin{kiconclusions}
	Tato práce popsala návrh a~implementaci kompilátoru zásobníkového jazyka
	inspirovaného jazykem \Forth{}. Výsledná implementace \\Roth{} v~jazyce Rust
	zahrnuje lexikální a~syntaktickou analýzu, sémantickou kontrolu, převod do
	mezikódu, optimalizační průchody a~generování cílového kódu.

	Součástí návrhu je zásobníkový mezikód a~sada optimalizačních průchodů
	(constant folding, inlining, peephole optimalizace, strength reduction
	a~eliminace mrtvého kódu). Implementace podporuje generování kódu do jazyků
	Rust a~C a~obsahuje REPL prostředí založené na průběžné kompilaci a dynamickém
	načítání knihoven.

	Další rozvoj může směřovat k~rozšíření backendů, doplnění dalších instrukcí
	mezikódu, rozšíření standardní knihovny a~k~systematičtějšímu měření výkonu
	na větší sadě referenčních programů.
\end{kiconclusions}

\begin{kiconclusions}[english]
	This thesis described the design and implementation of a stack-based language
	inspired by Forth. The resulting Rust implementation of Roth includes lexing,
	parsing, semantic validation, lowering to an intermediate representation,
	IR-level optimizations, and target code generation.

	The design uses a dedicated stack-based intermediate representation together
	with optimization passes such as constant folding, inlining, peephole
	optimization, strength reduction, and dead code elimination. The implementation
	supports Rust and C backends and includes a REPL environment based on
	incremental compilation and dynamic library loading.

	Future work may focus on extending backend coverage, adding more IR
	instructions, expanding the standard library, and performing more systematic
	performance measurements on a broader benchmark set.
\end{kiconclusions}
